% Wykres przepustowości dla różnych konfiguracji
% Wymagane pakiety w preambule:
% \usepackage{pgfplots}
% \pgfplotsset{compat=1.18}

\begin{figure}[h]
\centering
\begin{tikzpicture}
\begin{axis}[
    xlabel={Liczba cząstek (N)},
    ylabel={Przepustowość [Minteract/s]},
    legend pos=north west,
    grid=major,
    width=0.9\textwidth,
    height=0.6\textwidth,
    xmin=30, xmax=400,
    ymin=0, ymax=2800,
    xtick={50,100,150,200,250,300,350,400},
    legend style={font=\small},
    mark size=2pt
]

% n1c8 - czysty OpenMP
\addplot[color=blue,mark=o,thick] coordinates {
    (50, 339.20)
    (80, 524.01)
    (110, 229.79)
    (140, 238.14)
    (170, 303.21)
    (200, 202.74)
    (230, 252.19)
    (260, 404.15)
    (290, 321.19)
    (320, 284.77)
    (350, 302.96)
    (380, 371.19)
};
\addlegendentry{n1c8 (1 węzeł × 8 wątków)}

% n2c4 - hybrid
\addplot[color=green,mark=square,thick] coordinates {
    (50, 296.31)
    (80, 571.48)
    (110, 442.23)
    (140, 457.86)
    (170, 582.27)
    (200, 539.78)
    (230, 697.05)
    (260, 731.34)
    (290, 769.43)
    (320, 805.97)
    (350, 789.19)
    (380, 899.52)
};
\addlegendentry{n2c4 (2 węzły × 4 wątki)}

% n4c2 - hybrid
\addplot[color=orange,mark=triangle,thick] coordinates {
    (50, 116.10)
    (80, 245.33)
    (110, 430.94)
    (140, 739.15)
    (170, 902.97)
    (200, 1137.51)
    (230, 1323.93)
    (260, 1285.64)
    (290, 1537.42)
    (320, 1693.22)
    (350, 1740.91)
    (380, 1756.35)
};
\addlegendentry{n4c2 (4 węzły × 2 wątki)}

% n8c1 - czysty MPI
\addplot[color=red,mark=*,thick] coordinates {
    (50, 114.76)
    (80, 281.09)
    (110, 439.33)
    (140, 741.87)
    (170, 1013.86)
    (200, 1280.54)
    (230, 1431.50)
    (260, 1711.23)
    (290, 1949.13)
    (320, 1994.20)
    (350, 2364.33)
    (380, 2630.61)
};
\addlegendentry{n8c1 (8 węzłów × 1 wątek)}

\end{axis}
\end{tikzpicture}
\caption{Przepustowość symulacji N-body dla różnych konfiguracji równoległych (5000 iteracji)}
\label{fig:wydajnosc}
\end{figure}

% Wykres przyspieszenia (speedup) względem n1c8
\begin{figure}[h]
\centering
\begin{tikzpicture}
\begin{axis}[
    xlabel={Liczba cząstek (N)},
    ylabel={Przyspieszenie względem n1c8 [×]},
    legend pos=north west,
    grid=major,
    width=0.9\textwidth,
    height=0.6\textwidth,
    xmin=30, xmax=400,
    ymin=0, ymax=8,
    xtick={50,100,150,200,250,300,350,400},
    legend style={font=\small},
    mark size=2pt
]

% Linia bazowa (n1c8 = 1.0)
\addplot[color=blue,dashed,thick] coordinates {
    (50, 1.0) (380, 1.0)
};
\addlegendentry{n1c8 (baseline)}

% n2c4
\addplot[color=green,mark=square,thick] coordinates {
    (50, 0.87)
    (80, 1.09)
    (110, 1.92)
    (140, 1.92)
    (170, 1.92)
    (200, 2.66)
    (230, 2.76)
    (260, 1.81)
    (290, 2.40)
    (320, 2.83)
    (350, 2.61)
    (380, 2.42)
};
\addlegendentry{n2c4}

% n4c2
\addplot[color=orange,mark=triangle,thick] coordinates {
    (50, 0.34)
    (80, 0.47)
    (110, 1.88)
    (140, 3.10)
    (170, 2.98)
    (200, 5.61)
    (230, 5.25)
    (260, 3.18)
    (290, 4.79)
    (320, 5.95)
    (350, 5.75)
    (380, 4.73)
};
\addlegendentry{n4c2}

% n8c1
\addplot[color=red,mark=*,thick] coordinates {
    (50, 0.34)
    (80, 0.54)
    (110, 1.91)
    (140, 3.12)
    (170, 3.34)
    (200, 6.32)
    (230, 5.68)
    (260, 4.23)
    (290, 6.07)
    (320, 7.00)
    (350, 7.80)
    (380, 7.09)
};
\addlegendentry{n8c1}

\end{axis}
\end{tikzpicture}
\caption{Przyspieszenie różnych konfiguracji względem czystego OpenMP (n1c8)}
\label{fig:speedup}
\end{figure}
